% 2026-09-28：统计口径修订；历史实验结果不重写。
\chapter{概率、统计与配对实验}
\label{app:stats}

E01 的结果允许读者复算一个正向点估计，却不允许直接宣布持久更新更好。
本附录解释差值、区间、样本规模与研究结论之间的关系。
所有能力读数仍来自原始实验；样本量是根据试验方差计算的规划量。

\section{效应量与检验回答不同问题}

在指定零假设和统计模型下，p 值衡量观察到当前或更极端检验统计量的概率。
它不是零假设成立的概率，也不度量效应的工程重要性；
ASA 的声明明确区分了这些含义\cite{asa-p-values}。
样本量增加可能使很小的差异更容易被检测，但某一次样本序列的 p 值没有随题数严格单调的保证。

本书报告效应量和区间，并另列工程上值得关注的差异。
点估计小于工程门槛，只说明当前估计没有达到该门槛，不能当作精确零效应。
最小可检测效应描述检验设计中的目标差异；最小有意义效应描述工程决策。
二者数值可以相同，定义不能混用。

\section{配对设计与配对差}

两个方案在同一批题上运行，第 $i$ 题各自的重复运行均值为 $q_A(x_i)$ 与 $q_B(x_i)$，定义
\begin{equation}
 d_i=q_A(x_i)-q_B(x_i),\qquad
 \widehat\Delta=\frac1n\sum_{i=1}^{n}d_i.
 \label{eq:paired}
\end{equation}
NIST 的配对观察分析同样以对应测量之差为基本对象\cite{nist-paired}。
本书把对应关系落实到任务标识，避免比较来自不同题目的两个成绩。当前配对入口还检查每臂标识唯一、同题的族与划分一致，以及两侧判定非空且为二值；空记录须先说明缺失原因，无法直接计算通过率。

配对能否降方差取决于两臂结果的协方差。
同题难度使两臂正相关时，差值可以抵消部分共同波动；难度不会在任何情况下都被完全消去。
若某题缺一臂，可以计算完整配对上的描述结果，但须另报排除的题和原因。
困难题更易缺失时，完整配对结果可能偏向较易完成的题，不能无条件代表原计划题集。

\section{重复运行与重采样单位}

同一题的重复调用估计题内随机性，不能按调用次数增加独立任务数。
现有 E01 每题运行 \measured{5} 次，封存题数仍是 \measured{16}。
种子记录帮助追踪配置；只有后端实际接收并执行确定的采样规则时，才能进一步声称种子可复现。

题级和族级重采样对应不同的推广范围。
固定族内的新题推断可以以题为对象，但要检查题内和族内依赖；
对新任务族的主张需要足够多独立留出族，并让权重和区间对应族级设计。
族数不足时应限制结论，而不是退回题级区间后继续保留广泛跨族主张。

E01 的格 4 减格 1 在题级重采样下区间宽度为 \measured{0.275}，
族级为 \measured{0.300}；格 4 减格 2 的次序相反，分别为 \measured{0.413} 与 \measured{0.363}。
这些数来自该轮结果记录，仅描述当前样本。
两种区间宽度的次序还受族内协方差和加权方式影响，不能仅凭反序就诊断某一种估计失效。

\section{百分位自助区间}

现有实现位于 \texttt{experiments/lab/eval.py} 的 \texttt{paired\_diff}。
它按指定单位有放回重采样，每次计算配对均值，再用自助分布的端点分位数形成区间。
固定自助随机种子可以复算同一份数据的端点，不能让原始模型调用变为确定性运行。

历史实现的族级分支先抽族，再把抽中族的全部题合并求均值；族内题数不同时，这不等于后文教学例的等族平均。历史端点按排序数组中从零计数的$\lfloor0.025B\rfloor$与$\lfloor0.975B\rfloor$位置取值（$B$为正整数自助次数），与后文逆经验分布分位数的约定不同。复算时须同时保留统计量、权重与端点规则，不能仅凭“族级自助”这个名称替换实现。

这种方法不要求逐题通过率本身服从正态分布，但依然依赖样本代表性和合适的重采样单位。
少量族、离散边界、相关观测或自适应选择可能使覆盖率不理想。
增加自助抽样次数只能减少区间计算中的模拟误差，不能补充真实任务信息。
本书保留历史区间，同时将超出设计范围的推断列为未解决问题。

\section{区间精度与检验功效}

假设每道题给出一个配对差。希望均值的双侧区间足够窄，可以从目标半宽反推题数：
\begin{equation}
 n_{\mathrm{precision}}\approx
 \left\lceil\left(\frac{z_{1-\alpha/2}s}{h}\right)^2\right\rceil.
 \label{eq:samplesize}
\end{equation}
其中 $h$ 是目标半宽，$s$ 来自试验配对差，$\alpha$ 是双侧区间的错误概率。
在独立近似正态和方差估计可用的条件下，令 $h=0.03$ 并取 $z=1.96$，
E01 的历史计算得到 \calcnum{353} 与 \calcnum{836} 道配对题。
这里的 $0.03$ 限定均值估计的精度。

指定检验功效后，样本量还取决于漏检概率 $\beta$。
围绕零差异的双侧均值检验可先用
\begin{equation}
 n_{\mathrm{power}}\approx
 \left\lceil\left(
 \frac{(z_{1-\alpha/2}+z_{1-\beta})s}{\delta_{\mathrm{detect}}}
 \right)^2\right\rceil
\end{equation}
估算所需题数。NIST 在总体标准差已知时给出正态近似，标准差未知时还给出按自由度迭代的 $t$ 近似\cite{nist-sample-size}。这里用先导试验的 $s$ 代入正态公式，把它当作未来配对差波动大小的估计。

仍取双侧 $\alpha=0.05$。如果真实均值相差 $0.03$，希望有 $80\%$ 的机会检出它，E01 的两个比较分别需要约 \calcnum{721} 和 \calcnum{1708} 道配对题。以格 4 与格 1 为例，从 353 道增加到 721 道，原因在于功效公式还要为漏检留出余量。两个 $0.03$ 虽然数值相同，一个指定区间半宽，另一个指定要检出的真实差异。未舍入的先导标准差与两套计算均保存在 \path{analysis/sample-size-reproduction.json}。若按任务族抽样、同时比较多种策略或分批查看结果，则要按相应设计重新规划。

\section{从试验运行到确认性运行}

试验运行给出成本、方差与实现缺陷，用于规划后续采样。
下一轮在启动前冻结目标分布、主要比较、工程阈值、功效、分析单位、缺失处理与停止规则。
已经用于查看结果、修改方法或选择设计的数据保留为试验资料；新的确认性结果使用另行封存的数据。

原协议中的三个百分点目标保持不变。
本次修订纠正的是统计解释，不根据现有点估计调节目标以换取有利状态。
若工程需求后来改变，应记录新问题和新协议，并对新数据执行，不能把旧试验改称为已完成的新检验。

\section{怎样从统计结果形成结论}

统计结果需要按比较设计解释。完整运行、干预识别、任务覆盖和精度分别约束结论。表\ref{tab:verdicts}给出四种阅读情形；它们帮助判断证据的用途，不为研究问题分配统一状态。

\begin{table}[htbp]
\centering
\caption{统计结果的用途与形成结论所需的条件}
\label{tab:verdicts}
\begin{tabular}{p{2.0cm}p{10.5cm}}
\toprule
阅读情形 & 含义与必要条件 \\
\midrule
证据待分辨 & 设计、运行、样本、精度或判据尚不足以作所声称的判断；写明具体缺口。 \\
范围内支持 & 在预先规定的范围内得到符合该命题判据的支持；有限实验不证明无条件普遍命题。 \\
条件边界 & 有证据支持明确的条件边界；不能仅因一组区间跨零而进入。 \\
有效反证 & 有适合该命题的反证，且设计、方向和分析有效；不能由“没有检出”直接进入。 \\
\bottomrule
\end{tabular}
\end{table}

\begin{sloppypar}
区间跨零不证明精确相等，即使总题数达到某个规划值也一样。
条件性结论需要直接检验边界，等价结论需要预先定义可接受差异范围。
十个研究问题及比较设计见第16章；原假设、旧判据与状态记录见附录\ref{sec:appg-history}和\nolinkurl{spec/propositions.md}。
\end{sloppypar}

\section*{本附录用到的代码}
\begin{sloppypar}
\begin{itemize}
 \item \texttt{experiments/lab/eval.py}：配对与自助实现；当前入口增加输入检查，保留历史统计量、重采样和端点规则。返回的降级字段不直接决定命题状态。
 \item \texttt{calculations/audit\_sample\_size.py}：离线复算精度与功效示例，不调用模型、不生成能力收益。
 \item \nolinkurl{analysis/inputs.json}：按运行编号给出实际读取的汇总路径与 SHA-256。公开预览的 E03 条目另列原件的路径和摘要，可据此区分公开副本与维护档案原件。
\end{itemize}
\end{sloppypar}

\section{反馈一致性表应该怎样读}

第 4 章 E02 的历史字段 \texttt{mismatch\_rate} 比较成功指示与验收标签，不能解释为有害更新比例。
对可解析反馈，记不一致数为 $m$、条数为 $n$，比例为 $m/n$；还应另报无法解析的条数。
若用实验台验收结果直接构造成功指示，二者一致是实现关系，不是评审模型准确率。

归档中自评臂的 \measured{8} 条反馈对应 \measured{6} 次成功与 \measured{2} 次失败，
评审臂对应 \measured{5} 次成功与 \measured{3} 次失败。各自标签全对不等于所有输入上全对，
也不说明诊断文字、建议修复和未来收益正确。
若要估计错误类别上的能力，应分别报告成功与失败样本的判定，并说明候选样本的抽取方式。

标签、诊断、修复与迁移还可能有不同的观察单位。一个提交可以产生多份评审，
一条经验可以影响多个后续回合，一条任务流内又可能共享记忆。
增加同一提交的评审次数不等于增加独立提交；增加同题回合不等于增加任务族。
分析与规划要围绕所要推广的单位，而不是选择最容易得到大样本数的那一层。

历史 E02 中，中间模型反馈原文没有单独归档；标签与经验可以分析，诊断全文无法逐字复核。
后续再调用一次模型得到的是新观察，不能补作旧记录。缺失项应影响结论范围，而不应被推测填满。

\section{一份小表为什么会给出两个平均收益}
\label{sec:appa-weighted-example}

以下例题全部由本书构造，不来自E01或任何模型运行。目的是把“按题平均”和“按族平均”的区别算到底，而不是用新例题补充实验证据。设一个计划包含四个任务族，各族里新旧策略的平均得分差已经算好，差的正号表示新策略较好。为了突出权重，暂时把每个族的平均差当作已给定的数，不讨论它本身的估计误差。

\begin{table}[htbp]
\centering
\caption{权重教学例：输入为人为设定的题数和族内平均差}
\label{tab:appa-weight-example}
\begin{tabular}{lrrr}
\toprule
族 & 题数 & 族内平均差 & 题数乘平均差 \\
\midrule
A & \calcnum{8} & \calcnum{0.10} & \calcnum{0.80} \\
B & \calcnum{2} & \calcnum{-0.20} & \calcnum{-0.40} \\
C & \calcnum{2} & \calcnum{0.30} & \calcnum{0.60} \\
D & \calcnum{2} & \calcnum{-0.10} & \calcnum{-0.20} \\
\bottomrule
\end{tabular}
\end{table}

若每一道题获得相同权重，平均差为
\begin{equation}
 \widehat\Delta_{\mathrm{task}}
 =\frac{8(0.10)+2(-0.20)+2(0.30)+2(-0.10)}{14}
 =\frac{2}{35}\approx 0.05714.
\end{equation}
若每个任务族获得相同权重，则为
\begin{equation}
 \widehat\Delta_{\mathrm{family}}
 =\frac{0.10-0.20+0.30-0.10}{4}=0.025.
\end{equation}
两种计算都没有算术错误；前者对应从这份题表均匀抽题，后者对应先均匀抽族、再取该族的平均表现。A族在前一种计算里占\calcnum{8/14}，在后一种里占\calcnum{1/4}。由于A族的平均差高于等族平均，它获得较大权重时，总收益随之提高。由此得到的\calcnum{5.714}个百分点与\calcnum{2.5}个百分点回答不同的问题，不能只保留其中越过工程门槛的那个。

在实际研究中，题数可能取决于哪个任务生成器更容易产出可用样本，而不取决于真实部署频率。若直接把生成方便的族放大，研究者就把数据获取难度变成了目标人群权重。另一方面，业务系统确实可能经常遇到某一类问题，此时等族平均也未必是唯一合理目标。正确顺序是先写出要推广到的任务分布，再决定权重；不能看完两列结果才选择更好的一列。

这种差异还影响对命题的判断。P1若限定为某个业务流上的资源分配优势，可以使用事先确定的业务频率；若声称对一批平等关注的独立任务族都有普遍收益，就还需要族级结果及异质性说明。一个加权正平均不会抹去B、D两族的负差，更不能自行改写为“在所有族上有效”。应同时保留总体目标、各族表现和权重来源，让读者知道总体收益由谁贡献。

\subsection{先聚合哪一层，也是在定义观察对象}

一般地，记第$f$族第$i$题第$r$次重复的配对差为$d_{fir}$。一种等族、族内等题、题内等重复的统计量是
\begin{equation}
 \widehat\Delta
 =\frac{1}{F}\sum_{f=1}^{F}\frac{1}{n_f}
  \sum_{i=1}^{n_f}\frac{1}{R_{fi}}\sum_{r=1}^{R_{fi}}d_{fir}.
 \label{eq:appa-nested-mean}
\end{equation}
每个分母都有对象：$F$是族数，$n_f$是该族题数，$R_{fi}$是该题重复数，均要求为正整数。若某题因为故障而多重跑几次，把所有回合直接摊平成一列，会使它获得更大权重。式~\eqref{eq:appa-nested-mean}可以维持题权重，但不能消除“为什么这题被多跑”的选择问题；若重跑次数取决于已见成绩，还要处理停止规则带来的偏差。

还要区分“同一题的两臂配对”与“同一个随机数的两臂配对”。两臂做了相同题目，已经可以构成题目层面的对应；这不要求服务商保证同一个种子产生相同内部随机过程。没有共同随机数保证时，不要把采样序号相同描述为控制住全部随机性。保留题目、日期、后端标识、采样参数和重复标识，才有条件判断哪些变化得到控制，哪些只能进入误差和限制。

在线更新还增加了一层依赖。同一条任务流后面的题会读取前面形成的记忆，不能把它们任意打散后当作相互独立的新系统运行。若目标是比较完整任务流策略，统计记录应保留整条序列的起点、顺序和终点；若要比较固定版本，应先冻结版本，再执行新的评价。两种设计可以共享代码，但不能共享一个含义模糊的“样本数”。

\section{把四个族的自助分布完整列出来}

继续使用表~\ref{tab:appa-weight-example}的四个族平均差。这里只演示一种条件于已观察族均值的计算：每次有放回地抽四个族，每个位置等概率选择A、B、C、D，再计算等族平均。因为例子很小，全部有序抽样组合只有$4^4=\calcnum{256}$个，可以完整枚举，不必用随机数近似。

抽到$(A,A,A,A)$时，平均差为\calcnum{0.10}；抽到$(B,B,B,B)$时为\calcnum{-0.20}；抽到$(A,B,C,D)$时为\calcnum{0.025}。不同的有序组合可能得到同一个平均数，分布应保留这种重复次数。不能先去重再给每个不同平均值相同权重，否则改变了重采样概率。本例共有\calcnum{19}个不同平均值，但分位数仍按\calcnum{256}个等概率组合计算。

为使读者能够逐项复算，本书采用经验分布的逆分位数定义：将正整数$N$个值从小到大排列，对概率$0<p\le1$取一开始计数的第$\lceil pN\rceil$项，不作插值。由此，\calcnum{0.025}和\calcnum{0.975}分位数分别为\calcnum{-0.15}与\calcnum{0.20}。改变分位数插值约定可能改变某些小样本端点，因此代码、版本和分位数规则都是结果身份的一部分。

这里算出了一个自助百分位范围，却没有证明它对未知任务族总体有\calcnum{95\%}覆盖率。完整枚举只是消除了这一步的蒙特卡洛误差：给定这四个族均值和抽样规则，端点已无随机计算误差。它没有增加真实族数，没有估计族内重复的不确定性，也没有证明这四族具有代表性。尤其不能把\calcnum{256}种重组当作\calcnum{256}个独立任务族。

置信水平描述的是在相应假设下重复抽样并构造区间这一方法的覆盖性质，而不是在频率学派解释中为当前固定参数赋予“落进已得区间的概率”。NIST的均值置信区间说明明确作了这一区分\cite{nist-confidence-mean}。本例的用途是展示计算对象，不能借用该说明对均值区间的定义，为任意小样本自助规则自动提供覆盖保证。

若一个分析确实需要同时反映族、题和模型采样层面的变化，可以考虑与数据生成机制相匹配的多层重采样或模型。但“每一层都再抽一次”不是无需论证的通用配方。比如各族已经汇总成带有抽样误差的估计，再在不同层重复计入同一误差，可能改变所估计的方差；在线序列还必须保持时间上的共享状态。应先画出独立生成单位、条件依赖和目标量，再决定重采样哪些对象。

一个容易执行的审查办法，是给每个输入行加上族、题、重复、序列、版本五个标识，然后写一句话解释每次抽样会复制哪几列、保留哪几列。若抽中某个题，却把新旧策略的结果分别抽走，配对关系已经被破坏；若抽中某一序列中的末题，却不带它形成时的记忆来源，实际系统状态也被改了。自助法重组的是统计记录，不应悄悄造出实验中从未存在过的执行过程。

\section{为什么增加重复不能替代增加独立任务族}

下面用一个明确的教学模型展示方差来自哪里。假设配对差可写成
\begin{equation}
 d_{fir}=\Delta+u_f+v_{fi}+\epsilon_{fir},
\end{equation}
其中族偏移$u_f$、族内题偏移$v_{fi}$、重复噪声$\epsilon_{fir}$均为零均值，各层之间以及同层独立单位之间互不相关；方差分别为$\tau_F^2$、$\tau_T^2$和$\sigma^2$。每族都有$T$题、每题都有$R$次重复时，等权总平均的方差由方差加法直接得到
\begin{equation}
 \operatorname{Var}(\widehat\Delta)
 =\frac{\tau_F^2}{F}+\frac{\tau_T^2}{FT}
   +\frac{\sigma^2}{FTR}.
 \label{eq:appa-variance-levels}
\end{equation}
这是本例在所列独立性与平衡设计假设下的代数关系，不是对E01二元结果拟合出的模型。真实更新策略可能存在更复杂的协方差、不同题的重复方差和时间漂移，不能直接套用本例参数规划付费试验。

为便于看量级，人为指定$\tau_F^2=\calcnum{0.04}$、$\tau_T^2=\calcnum{0.09}$、$\sigma^2=\calcnum{0.16}$。当$F=\calcnum{4}$、$T=\calcnum{5}$、$R=\calcnum{2}$时，总平均方差为\calcnum{0.0185}，标准误约\calcnum{0.1360}。只把每题重复数提高到\calcnum{20}，方差变为\calcnum{0.0149}，标准误约\calcnum{0.1221}；重复增加很多，族偏移与题偏移两项仍在。

若保持每族\calcnum{5}题、每题\calcnum{2}次重复，而把独立族数从\calcnum{4}提高到\calcnum{40}，则在这一教学模型下方差变为\calcnum{0.00185}，标准误约\calcnum{0.0430}。这不是建议任何项目直接采集四十族，也没有比较获取新族与重复调用的费用；它只说明，在族间差异主导时，继续压低题内噪声存在结构上的收益上限。

实验规划因而至少要分开两种成本：开发一个具有独立语义和失效机制的新族的成本，以及在已有题上再调用一次模型的成本。新族往往更难构造和验证，不能仅按API单价评价；廉价重跑则未必增加研究最缺的那类信息。小试可以帮助了解这些量，却不能把小试方差当作稳定已知的自然常数。规划时应报告多个合理参数下的敏感性，而不是只挑一个能够使预算最小的点估计。

式~\eqref{eq:appa-variance-levels}还给出一个读表顺序：先确认目标是新回合、新题还是新族，再找相应的不可被重复抵消的项。若只关心当前固定题集在未来模型采样中的平均表现，可以条件于这批题讨论重复噪声；若论文标题声称跨域能力，便不能在分析中把族差异条件掉，再继续使用同样宽泛的标题。统计单位和结论范围必须一起变化。

\section{缺失配对的保守范围怎样计算}

另设一个独立教学场景：计划比较\calcnum{14}道题，只取得\calcnum{10}道完整配对，它们的差值总和为\calcnum{1}。完整配对均值因此是\calcnum{0.10}。对于成功率差，每个尚未观察的配对差只能确定在$[-1,1]$内。若不额外假设缺失机制，全部计划题平均差的逻辑范围是
\begin{equation}
 \frac{1-4}{14}\le\Delta_{\mathrm{planned}}
 \le\frac{1+4}{14},
\end{equation}
即约$[\calcnum{-0.2143},\calcnum{0.3571}]$。这个范围不是置信区间，没有覆盖水平；它描述在所给取值边界内，未观察差值可能把完整计划的平均推到哪里。它与完整配对上的正均值可以同时成立，因此“已完成题的平均提高”并不自动代表计划题集的平均提高。

这个粗界只用了完整对数与总和。若缺失题的一臂结果已知，界可以更窄：已知旧策略成功而新策略缺失时，差在$[-1,0]$；已知旧策略失败而新策略缺失时，差在$[0,1]$。若一题有部分重复，也可按事先规定的题内分母分解已知和未知部分。信息越细，边界可以越具体，但每一步都要说明使用了什么观察，不能把未知结果按倾向填成成功。

失败与缺失还要按产生机制区分。若策略在既定预算内未提交任何合法候选，这是完整策略表现的一部分，可以按预定规则记失败；若服务中断导致本应执行的比较完全没有响应，它是外部观察缺失，需另记范围与敏感性。若只对更新臂的网络错误重跑，而对重试臂按失败计分，两臂执行条件就不再相同。处理规则应在看见方向之前确定，并保留中断记录与全部已发生费用。

不要用“我们做了敏感性分析”代替具体假设。至少应列出基准分母、缺失行、保守边界以及一种有业务理由的替代情形。若替代情形依赖“缺失题与完成题难度相同”，就写出这一假设的依据与不可验证部分。敏感性结果的作用是显示结论依赖什么，不是从多个设定中再选一个支持原主张的结果。

\section{单位资源收益的两种平均也不一样}

再考虑两条人为构造的序列：收益分别为\calcnum{0.10}与\calcnum{0.20}，成本分别为\calcnum{2}与\calcnum{10}个同种资源单位。先对每条序列求收益成本比，再等权平均，得到
\begin{equation}
 \frac12\left(\frac{0.10}{2}+\frac{0.20}{10}\right)=0.035.
\end{equation}
把收益相加后除以总成本，则为
\begin{equation}
 \frac{0.10+0.20}{2+10}=0.025.
\end{equation}
前者把两条序列的效率各看作一个观察，后者按消耗资源加权。若研究问题是固定预算可以获得多少累计收益，后者通常更贴近所写的资源分配目标；若研究对象是随机抽到一条序列的效率，则需要明确该序列分布以及成本接近零时如何处理。两者不应共享一个没有分母说明的“平均效率”标签。

收益成本比还有一个基础前提：分子必须能够按所用方式相加。本例只作代数展示；真实准确率是比例，跨序列相加时要先明确题数和权重。如果序列长度不同，直接把两个准确率增益相加并不等于增加的正确任务总数。可分别报告预定平均增益、总成功任务增量和总资源，避免一个比值同时隐藏分子、分母和权重三种选择。

成本中若含未知费用，不能让缺失自然变成零。已知token和未知计费用量应并列展示；若美元价格没有可靠版本，只报告实物资源。共享准备费在研究实际开销中计一次，在假设分别部署每个方案时则要为各方案说明取得该准备物的费用。把后一个反事实成本表再次求和，会重复计算实际支出；把前一个共享账当作每个独立方案都免费得到记忆，又会低估部署成本。

\section{五对观测的单侧检验与双侧检验}
\label{sec:appa-signed-rank}

小样本的显著性判断还取决于检验方向。下面完全采用构造数据，说明为什么“只有五次运行，所以$p<0.05$必然有误”并不成立。设预先定义的五个独立配对差为$\calcnum{(1,2,3,4,5)}$，绝对值没有并列，也没有零差。Wilcoxon符号秩检验先按绝对值排序，再计算正差的秩和$T_+$。这里五个差都为正，故$T_+=\calcnum{15}$。

在差值来自以零为中心的对称分布这一零假设下，给定这些绝对值，五个符号的$2^5=\calcnum{32}$种组合等概率。只有全正这一种组合使$T_+$达到最大值；若事先选择正向单侧备择，则精确尾部概率为
\[
 p_{\mathrm{one}}=\frac{1}{32}=\calcnum{0.03125}.
\]
若事先选择双侧检验，另一端全负同样极端，本例按两倍较小尾部的定义得到
\[
 p_{\mathrm{two}}=\frac{2}{32}=\calcnum{0.0625}.
\]
同一组配对差因此可以在单侧规则下低于\calcnum{0.05}，在双侧规则下高于该值。不能看完方向再把双侧改成有利的单侧，也不能从一句“小于阈值”反推出作者采用了哪个规则。

SciPy的版本化文档分别列出备择方向、零差处理和尾部计算方法，并提示并列秩、零差和浮点减法可能影响结果\cite{scipy-wilcoxon-116}。本例完整枚举的是没有零差或并列绝对值时的条件符号分布；有零差、并列或依赖时，不能只把运行数代入$2^n$照搬上面的结论。配对关系本身也要有依据：同一行号不自动表示相同任务起点或匹配的随机条件。

这正是阅读第12章迁移统计时需要补齐的信息。论文报告五次运行和Wilcoxon检验，可以与单侧显著性相容；只有绑定实际比较的逐运行记录、配对、方向、零差处理及软件版本，才有条件复核原报告值。这个教学枚举解释了可能的统计机制，没有重建来源论文的实验，也没有为缺失设置代填默认值。

\section{从统计例题回到研究问题}

以上六组例题分别改变权重、重采样对象、采样层级、观察完整性、资源比值的聚合方式和检验方向，没有一组产生本书命题的新结果。它们要求实验报告在效应量之前先交代目标量的组成：哪些任务属于目标，如何加权，哪一层独立，哪一部分费用已知，以及哪些输出未被观察。缺少其中任何一项，都可能使一个精确到小数点后四位的数字回答了另一个问题。

对于E01新版，这意味着同源记忆、重试选择和调用总账修好以后，仍要另外完成目标族、样本规模和分析层级。软件检查能证明选择器没有读取其接口之外的字段，不能证明族样本具有代表性；保存了所有调用，也不能证明预算对业务目标合理。实现修复、设计有效性和统计精度是连续但不可互换的验收步骤。

本附录的确定性复算位于\texttt{calculations/reproduce\_appendix\_a.py}，输出包含输入表、分位数定义、完整重组分布以及所有公式的参考结果。可从仓库根目录运行
\begin{lstlisting}[language=bash]
python calculations/reproduce_appendix_a.py -o reader-work/appendix-a.json
\end{lstlisting}
将当前完整结果另存到新路径；目标已存在时拒绝覆盖。省略\texttt{-o}选项时沿用\path{analysis/appendix-a-reproduction.json}默认位置并更新该派生报告。脚本不导入模型后端、不生成实验样本。读者可以改变教学输入观察量的变化，但改变后的数值仍属于新教学设定；除非真的执行了符合新协议的实验，它们不会变成新的模型证据。
